%% toc: P2 20 20.7,20.8
\section{Proofs on inclusion form, nonstandard models and countability}
\label{sec:annex-p02-w02-p02-033-inclusion-and-skolem}

Notation is that of \cref{sec:p02-w02-p02-033-inclusion-and-skolem}.  For a
tuple \(\bar a=(a_1,\ldots,a_k)\) and a map \(h\), \(h\bar a\) is
\((h(a_1),\ldots,h(a_k))\).

\begin{annexitem}{annex:p02-w02-p02-033-embed}{thm:p02-w02-p02-033-elementary-maps-embed}{Elementary embeddings are embeddings}
\begin{proof}
Apply elementarity to four kinds of formula.
\emph{Injectivity.} If \(a\ne b\) in \(\dom\mathcal M\), then
\(\mathcal M\entails\neg(x_1=x_2)\) at \((a,b)\), so
\(\mathcal N\entails\neg(x_1=x_2)\) at \((h(a),h(b))\), that is,
\(h(a)\ne h(b)\).
\emph{Constants.} For a constant symbol \(c\), the formula \(x_1=c\) holds in
\(\mathcal M\) at \(c^{\mathcal M}\); hence it holds in \(\mathcal N\) at
\(h(c^{\mathcal M})\), so \(h(c^{\mathcal M})=c^{\mathcal N}\).
\emph{Functions.} For an \(n\)-ary function symbol \(f\) and
\(\bar a\in(\dom\mathcal M)^n\), the formula
\(f(x_1,\ldots,x_n)=x_{n+1}\) holds in \(\mathcal M\) at
\((\bar a,f^{\mathcal M}(\bar a))\); hence
\(f^{\mathcal N}(h\bar a)=h(f^{\mathcal M}(\bar a))\).
\emph{Relations.} For a relation symbol \(R\), elementarity applied to
\(R(x_1,\ldots,x_n)\) gives
\(\mathcal M\entails R(\bar a)\) iff \(\mathcal N\entails R(h\bar a)\).
So \(h\) is an injective homomorphism that reflects every relation symbol,
an embedding by the definition in \cref{sec:maps-elementarity-01}.  A
sentence is a formula with no free variable, and elementarity at the empty
tuple gives \(\mathcal M\entails\varphi\) iff \(\mathcal N\entails\varphi\).

If \(\dom\mathcal M\subseteq\dom\mathcal N\) and \(h\) is the inclusion
map, the inclusion is therefore an embedding, which is the definition of
\(\mathcal M\subseteq\mathcal N\) in \cref{sec:maps-elementarity-02}; and it
is elementary, which is the definition of
\(\mathcal M\preccurlyeq\mathcal N\) in \cref{sec:maps-elementarity-03}.
\end{proof}
\end{annexitem}

\begin{annexitem}{annex:p02-w02-p02-033-copy}{thm:p02-w02-p02-033-copy-over-image}{Isomorphic copy over the image of an elementary embedding}
\begin{proof}
Write \(M=\dom\mathcal M\) and \(N=\dom\mathcal N\), so that
\(D=N\setminus j[M]\).  The elements of \(M\) and \(N\) are sets of the
ambient set theory, and \(\langle u,v\rangle=\{\{u\},\{u,v\}\}\) is the
Kuratowski pair.

\emph{Step 1: a fresh index.}  Let
\[
 R=\Bigl\{u\in\textstyle\bigcup\bigcup M:
 \langle u,d\rangle\in M\text{ for some }d\in D\Bigr\},
\]
a set by union and separation, and let \(r=\{u\in R:u\notin u\}\).  If
\(r\in R\), then \(r\in r\) iff (\(r\in R\) and \(r\notin r\)) iff
\(r\notin r\), which is impossible; so \(r\notin R\).  If
\(\langle u,d\rangle\in M\) with \(d\in D\), then \(u\in\{u\}\in
\langle u,d\rangle\in M\), so \(u\in\bigcup\bigcup M\) and \(u\in R\).
Hence \(\langle r,d\rangle\notin M\) for every \(d\in D\).  Put
\(D'=\{\langle r,d\rangle:d\in D\}\); then \(M\cap D'=\varnothing\).

\emph{Step 2: the bijection.}  Let \(N'=M\cup D'\) and define
\(\pi:N'\to N\) by \(\pi(a)=j(a)\) for \(a\in M\) and
\(\pi(\langle r,d\rangle)=d\) for \(d\in D\).  This is well defined because
\(M\) and \(D'\) are disjoint.  It is injective: \(j\) is injective by
\cref{thm:p02-w02-p02-033-elementary-maps-embed}; \(\langle r,d\rangle=
\langle r,d'\rangle\) forces \(d=d'\); and \(\pi[M]=j[M]\) and
\(\pi[D']=D\) are disjoint.  It is surjective because
\(j[M]\cup D=N\).

\emph{Step 3: transport of structure.}  Define the \(\Lang\)-structure
\(\mathcal N'\) with domain \(N'\) by
\[
\begin{aligned}
 c^{\mathcal N'}&=\pi^{-1}(c^{\mathcal N}),\\
 f^{\mathcal N'}(\bar x)&=\pi^{-1}\bigl(f^{\mathcal N}(\pi\bar x)\bigr),\\
 R^{\mathcal N'}&=\{\bar x\in(N')^n:\pi\bar x\in R^{\mathcal N}\}
\end{aligned}
\]
for constant symbols \(c\), \(n\)-ary function symbols \(f\) and \(n\)-ary
relation symbols \(R\).  Applying \(\pi\) gives
\(\pi(c^{\mathcal N'})=c^{\mathcal N}\) and
\(\pi(f^{\mathcal N'}(\bar x))=f^{\mathcal N}(\pi\bar x)\), and
\(\bar x\in R^{\mathcal N'}\) iff \(\pi\bar x\in R^{\mathcal N}\).  So
\(\pi\) is a homomorphism that reflects every relation symbol; being
bijective, it is a surjective embedding, that is, an isomorphism
\(\mathcal N'\to\mathcal N\).

\emph{Step 4: elementarity of the inclusion (clause 1).}  Let
\(\varphi(x_1,\ldots,x_k)\) be an \(\Lang\)-formula and
\(\bar a\in M^k\).  By \cref{lem:iso-invariance} applied to \(\pi\),
\(\mathcal N'\entails\varphi(\bar a)\) iff
\(\mathcal N\entails\varphi(\pi\bar a)\).  Since \(\pi\bar a=j\bar a\) and
\(j\) is elementary, this holds iff \(\mathcal M\entails\varphi(\bar a)\).
Thus the inclusion map \(M\to N'\) is an elementary embedding
\(\mathcal M\to\mathcal N'\), and
\cref{thm:p02-w02-p02-033-elementary-maps-embed} gives
\(\mathcal M\subseteq\mathcal N'\) and \(\mathcal M\preccurlyeq\mathcal N'\).

\emph{Step 5: clause 2.}  Let \(\mathcal N'\) and \(\pi\) be any pair as in
clause~1.  Since \(\pi\) is a bijection, \(\abs{\dom\mathcal N'}=\abs N\).
Its restriction to \(M\) is \(j\), with image \(j[M]\); as \(\pi\) is a
bijection of \(\dom\mathcal N'\) onto \(N\), it maps the complement
\(\dom\mathcal N'\setminus M\) bijectively onto the complement
\(N\setminus j[M]=D\).  Hence \(\dom\mathcal N'\ne M\) iff
\(D\ne\varnothing\), that is, iff the range of \(j\) is proper.

\emph{Step 6: clause 3.}  First, the inverse \(\pi_2^{-1}\) of an
isomorphism is an isomorphism: it is a bijection, and applying
\(\pi_2^{-1}\) to the identities
\(\pi_2(c^{\mathcal N'_2})=c^{\mathcal N}\) and
\(\pi_2(f^{\mathcal N'_2}(\bar y))=f^{\mathcal N}(\pi_2\bar y)\) at
\(\bar y=\pi_2^{-1}\bar z\) gives
\(\pi_2^{-1}(c^{\mathcal N})=c^{\mathcal N'_2}\) and
\(\pi_2^{-1}(f^{\mathcal N}(\bar z))=f^{\mathcal N'_2}(\pi_2^{-1}\bar z)\),
while \(\bar z\in R^{\mathcal N}\) iff
\(\pi_2^{-1}\bar z\in R^{\mathcal N'_2}\) because \(\pi_2\) preserves and
reflects \(R\).  Second, a composite of isomorphisms is an isomorphism: it
is a bijection, and each defining clause for the composite follows by
applying the clause for each factor in turn.  So
\(\pi_2^{-1}\circ\pi_1:\mathcal N'_1\to\mathcal N'_2\) is an isomorphism,
and for \(a\in M\) it sends \(a\) to \(\pi_2^{-1}(j(a))=a\), because
\(\pi_2(a)=j(a)\).

No step selects an element from a family of sets without a rule: \(r\),
\(D'\), \(\pi\) and the interpretations are each defined outright, so the
Axiom of Choice is not used.
\end{proof}
\end{annexitem}

\begin{annexitem}{annex:p02-w02-p02-033-lsincl}{thm:p02-w02-p02-033-upward-ls-inclusion}{Upward L\"owenheim--Skolem in inclusion form}
\begin{proof}
Under the stated hypotheses,
\cref{thm:p02-upward-ls-elementary-extension} gives an
\(\Lang\)-structure \(\mathcal N\) with \(\abs{\dom\mathcal N}=\kappa\) and
an elementary embedding \(j:\mathcal M\to\mathcal N\) whose range is
proper.  Clause~1 of \cref{thm:p02-w02-p02-033-copy-over-image} gives
\(\mathcal N'\) with \(\mathcal M\preccurlyeq\mathcal N'\), and clause~2
gives \(\abs{\dom\mathcal N'}=\kappa\) and \(\dom\mathcal N'\ne\dom\mathcal
M\).

For the equivalent formulation, let \(\mathcal N'_{M}\) be the
\(\Lang(\dom\mathcal M)\)-expansion of \(\mathcal N'\) with
\(c_a^{\mathcal N'_M}=a\).  For a sentence
\(\varphi(c_{a_1},\ldots,c_{a_k})\) of \(\Lang(\dom\mathcal M)\), obtained
from an \(\Lang\)-formula \(\varphi(x_1,\ldots,x_k)\), the substitution
lemma \cref{thm:fol-substitution}, applied once for each closed term
\(c_{a_i}\), whose value is \(a_i\), together with
\cref{lem:reduct-invariance} give
\(\mathcal N'_M\entails\varphi(c_{a_1},\ldots,c_{a_k})\) iff
\(\mathcal N'\entails\varphi(a_1,\ldots,a_k)\), and the same holds for the
canonical expansion of \(\mathcal M\).  Hence \(\mathcal N'_M\) satisfies
every sentence of \(\operatorname{ElDiag}(\mathcal M)\) iff the inclusion
map is elementary, which by
\cref{thm:p02-w02-p02-033-elementary-maps-embed} and the definition of
\(\preccurlyeq\) is \(\mathcal M\preccurlyeq\mathcal N'\) once
\(\dom\mathcal M\subseteq\dom\mathcal N'\).  This is the inclusion-map case
of clause~2 of \cref{thm:diagram-tests}.
\end{proof}
\end{annexitem}

\begin{annexitem}{annex:p02-w02-p02-033-pure}{thm:p02-w02-p02-033-pure-sets}{Elementarity in the empty signature}
\begin{proof}
\emph{Clause 1.}  In the empty signature every injection preserves and
reflects all atomic formulas, since they are equalities; so the inclusion
\(\dom\mathcal M\to\dom\mathcal N\) is an embedding and
\(\mathcal M\subseteq\mathcal N\).  We verify the witness condition of
\cref{thm:tarski-vaught}.  Let \(\bar a=(a_1,\ldots,a_k)\) be a tuple from
\(\dom\mathcal M\) and suppose
\(\mathcal N\entails\exists y\,\varphi(y,\bar a)\), with witness \(b\in
\dom\mathcal N\).  If \(b\in\dom\mathcal M\) there is nothing to prove.
Otherwise, since \(\dom\mathcal M\) is infinite, choose
\(b'\in\dom\mathcal M\setminus\{a_1,\ldots,a_k\}\).  Let \(\sigma\) be the
permutation of \(\dom\mathcal N\) exchanging \(b\) and \(b'\) and fixing
every other element.  It fixes each \(a_i\), because \(a_i\ne b'\) by choice
and \(a_i\ne b\) as \(a_i\in\dom\mathcal M\) and \(b\notin\dom\mathcal M\).
In the empty signature a bijection is an isomorphism, so
\cref{lem:iso-invariance} gives
\(\mathcal N\entails\varphi(\sigma b,\sigma\bar a)\), that is,
\(\mathcal N\entails\varphi(b',\bar a)\) with \(b'\in\dom\mathcal M\).  The
Tarski--Vaught test gives \(\mathcal M\preccurlyeq\mathcal N\).

\emph{Clause 2.}  Let \(E=j[\dom\mathcal M]\) and let \(\mathcal E\) be the
set \(E\) as a structure.  The set \(E\) is infinite because \(j\) is
injective, so clause~1 gives \(\mathcal E\preccurlyeq\mathcal N\).  The map
\(j\) with codomain \(E\) is a bijection, hence an isomorphism
\(\mathcal M\to\mathcal E\).  For every formula \(\varphi\) and tuple
\(\bar a\) from \(\dom\mathcal M\), \cref{lem:iso-invariance} gives
\(\mathcal M\entails\varphi(\bar a)\) iff
\(\mathcal E\entails\varphi(j\bar a)\), and
\(\mathcal E\preccurlyeq\mathcal N\) gives this iff
\(\mathcal N\entails\varphi(j\bar a)\).  So \(j\) is elementary.
\end{proof}
\end{annexitem}

\begin{annexitem}{annex:p02-w02-p02-033-hyp}{thm:p02-w02-p02-033-hypotheses-needed}{No hypothesis of upward L\"owenheim--Skolem can be dropped}
\begin{proof}
\emph{Clause 1.}  Let \(\dom\mathcal M\) have exactly \(n\) elements, where
\(n\geq1\).  The sentence
\[
 \exists x_1\cdots\exists x_n\Bigl(\bigwedge_{1\leq i<i'\leq n}\neg(x_i=x_{i'})
 \land\forall y\,\bigvee_{1\leq i\leq n}y=x_i\Bigr)
\]
is true in \(\mathcal M\), so by
\cref{thm:p02-w02-p02-033-elementary-maps-embed} it is true in
\(\mathcal N\), and \(\dom\mathcal N\) has exactly \(n\) elements.  The map
\(j\) is injective by the same lemma, so its range has \(n\) elements and is
all of \(\dom\mathcal N\).  An inclusion-form extension gives such a \(j\)
(the inclusion map), so it is not proper either.

\emph{Clause 2.}  An elementary embedding is injective
(\cref{thm:p02-w02-p02-033-elementary-maps-embed}), so its target has
cardinality at least \(\abs{\dom\mathcal M}\).

\emph{Clause 3.}  The map \(\operatorname{code}\) is injective on finite
\(0\)--\(1\) sequences: for \(s\) of length \(\ell\), the number
\(\operatorname{code}(s)\) satisfies
\(2^{\ell}\leq\operatorname{code}(s)<2^{\ell+1}\), so it determines
\(\ell\), and then \(\operatorname{code}(s)-2^{\ell}=\sum_{i<\ell}s_i2^i\)
determines each \(s_i\) by uniqueness of binary expansion.
\(\mathcal M\) is infinite because its domain is \(\N\).  The map
\(x\mapsto f_x\) is injective and \(x\mapsto\{n:x(n)=1\}\) is a bijection
\(\{0,1\}^{\N}\to\mathcal P(\N)\), so
\(\abs{\Lang}\geq\abs{\mathcal P(\N)}>\aleph_0\) by Cantor's theorem in
\cref{sec:set-prerequisites-05}.

Let \(j:\mathcal M\to\mathcal N\) be an elementary embedding and
\(b\in\dom\mathcal N\setminus j[\N]\).  For \(x\ne y\) in
\(\{0,1\}^{\N}\), let \(m\) be least with \(x(m)\ne y(m)\).  For every
\(n>m\) the sequences \((x(0),\ldots,x(n-1))\) and \((y(0),\ldots,y(n-1))\)
differ at place \(m\), so injectivity of \(\operatorname{code}\) gives
\(f_x^{\mathcal M}(n)\ne f_y^{\mathcal M}(n)\).  Hence \(\mathcal M\)
satisfies the sentence
\[
 \forall z\,\Bigl(f_x(z)=f_y(z)\to\bigvee_{n\leq m}z=c_n\Bigr).
\]
By \cref{thm:p02-w02-p02-033-elementary-maps-embed}, \(\mathcal N\)
satisfies it as well, and \(c_n^{\mathcal N}=j(c_n^{\mathcal M})=j(n)\).
Since \(b\ne j(n)\) for every \(n\), the antecedent fails at \(z=b\):
\(f_x^{\mathcal N}(b)\ne f_y^{\mathcal N}(b)\).  Thus
\(x\mapsto f_x^{\mathcal N}(b)\) is an injection
\(\{0,1\}^{\N}\to\dom\mathcal N\), and
\(\abs{\dom\mathcal N}\geq\abs{\mathcal P(\N)}\).  For \(\kappa=\aleph_0\) a
target of cardinality \(\kappa\) would give an injection of
\(\mathcal P(\N)\) into a countable set, contradicting Cantor's theorem.  An
inclusion-form proper extension of cardinality \(\aleph_0\) would give such
a target with the inclusion map as \(j\).
\end{proof}
\end{annexitem}

\begin{annexitem}{annex:p02-w02-p02-033-numeral}{thm:p02-w02-p02-033-numeral-map}{The numeral map of a model of true arithmetic}
\begin{proof}
Write \(\operatorname{num}\) for \(\operatorname{num}_{\mathcal M}\).

\emph{Step 0: values of numerals.}  For every \(\Lang_A\)-structure
\(\mathcal K\), the term-evaluation clauses and \(\bar n=S^n0\) give
\(\bar 0^{\mathcal K}=0^{\mathcal K}\) and
\(\overline{n+1}^{\mathcal K}=S^{\mathcal K}(\bar n^{\mathcal K})\).  In
\(\mathcal N_{\mathrm{ar}}\), induction on \(n\) gives
\(\bar n^{\mathcal N_{\mathrm{ar}}}=n\).  In \(\mathcal M\) the same clauses
give \(\operatorname{num}(0)=0^{\mathcal M}\) and
\(\operatorname{num}(n+1)=S^{\mathcal M}(\operatorname{num}(n))\).

\emph{Step 1: substituting numerals.}  Let \(\varphi\) have its free
variables among \(x_1,\ldots,x_k\) and let \(n_1,\ldots,n_k\in\N\).  For
every \(\Lang_A\)-structure \(\mathcal K\),
\[
 \mathcal K\entails\varphi(\bar n_1,\ldots,\bar n_k)
 \quad\Longleftrightarrow\quad
 \mathcal K\entails\varphi(\bar n_1^{\mathcal K},\ldots,\bar n_k^{\mathcal K}).
\]
This is \(k\) applications of \cref{thm:fol-substitution}, substituting for
\(x_1\), then \(x_2\), and so on: numerals contain no variables, so the
value of each is independent of the assignment and later substitutions do
not act inside earlier ones.

\emph{Clause 1.}  Let a formula \(\varphi(x_1,\ldots,x_k)\) and numbers
\(n_1,\ldots,n_k\) be given, and put
\(\sigma=\varphi(\bar n_1,\ldots,\bar n_k)\).  By Steps~0 and~1,
\[
 \mathcal N_{\mathrm{ar}}\entails\varphi(n_1,\ldots,n_k)
 \iff\mathcal N_{\mathrm{ar}}\entails\sigma
 \iff\sigma\in\operatorname{Th}(\mathcal N_{\mathrm{ar}}).
\]
If
\(\sigma\in\operatorname{Th}(\mathcal N_{\mathrm{ar}})\), then
\(\mathcal M\entails\sigma\).  If not, then
\(\neg\sigma\in\operatorname{Th}(\mathcal N_{\mathrm{ar}})\), so
\(\mathcal M\entails\neg\sigma\) and \(\mathcal M\not\entails\sigma\).
Hence \(\sigma\in\operatorname{Th}(\mathcal N_{\mathrm{ar}})\) iff
\(\mathcal M\entails\sigma\), which by Step~1 holds iff
\[
 \mathcal M\entails\varphi\bigl(\operatorname{num}(n_1),\ldots,
 \operatorname{num}(n_k)\bigr).
\]
So \(\operatorname{num}\) is elementary.

\emph{Clause 2.}  First, \(\theta_n\) and \(\tau_n\) are true in
\(\mathcal N_{\mathrm{ar}}\).  For \(x\in\N\),
\(\mathcal N_{\mathrm{ar}}\entails x<\bar n\) iff there is \(z\in\N\) with
\(x+(z+1)=n\), iff \(x<n\) in the usual order, iff \(x\) is the value of
\(\bar k\) for some \(k<n\) (Step~0); for \(n=0\) both sides are false,
matching the empty disjunction.  For \(\tau_n\): if \(x\leq n\), then \(x\) is
the value of \(\bar x\) with \(x\leq n\); if \(x>n\), then
\(z=x-n-1\) satisfies \(n+(z+1)=x\), so
\(\mathcal N_{\mathrm{ar}}\entails\bar n<x\).  Hence
\(\theta_n,\tau_n\in\operatorname{Th}(\mathcal N_{\mathrm{ar}})\), and
\(\mathcal M\) satisfies both.  Applying \(\theta_n\) at \(b\) gives
\(\mathcal M\entails b<\bar n\) iff \(b=\bar k^{\mathcal M}\) for some
\(k<n\).  If \(b\) is nonstandard, then \(b\ne\bar k^{\mathcal M}\) for every
\(k\leq n\), and \(\tau_n\) at \(b\) leaves \(\mathcal M\entails\bar n<b\).

\emph{Clause 3.}  If \(\operatorname{num}\) is surjective, it is an
elementary embedding by clause~1, hence an embedding by
\cref{thm:p02-w02-p02-033-elementary-maps-embed}, and a surjective embedding
is an isomorphism by definition (\cref{sec:maps-elementarity-01}).
Conversely, let \(\mathcal M\cong\mathcal N_{\mathrm{ar}}\).  The inverse of
an isomorphism is a bijection that preserves constants and functions and
preserves and reflects relations, because the isomorphism reflects and
preserves them; so there is an isomorphism
\(h:\mathcal N_{\mathrm{ar}}\to\mathcal M\).  By
\cref{lem:term-preservation} and Step~0,
\(h(n)=h(\bar n^{\mathcal N_{\mathrm{ar}}})=\bar n^{\mathcal M}=
\operatorname{num}(n)\), so \(\operatorname{num}=h\) is surjective.  The
last equivalence is the definition of a standard element.

\emph{Clause 4.}  The following sentences are true in
\(\mathcal N_{\mathrm{ar}}\), so \(\mathcal M\) satisfies them:
\[
\begin{aligned}
 \alpha_1&=\forall x\,(\neg\,x=0\to\exists y\,(Sy=x)),&
 \alpha_2&=\forall x\,\forall y\,(Sx=Sy\to x=y),\\
 \alpha_3&=\forall y\,(y<Sy),&
 \alpha_4&=\forall x\,\neg(x<x),\\
 \alpha_5&=\forall x\,\forall y\,\forall w\,
   (x<y\land y<w\to x<w).
\end{aligned}
\]
In \(\mathcal N_{\mathrm{ar}}\): every nonzero number is a successor; \(S\)
is injective; \(y+(0+1)=y+1\) witnesses \(\alpha_3\); \(x+(z+1)=x\) has no
solution; and \(x+(z+1)=y\), \(y+(z'+1)=w\) give
\(x+((z+z'+1)+1)=w\).

Let \(b\) be nonstandard.  Then \(b\ne\operatorname{num}(0)=0^{\mathcal M}\),
so \(\alpha_1\) gives \(b'\) with \(S^{\mathcal M}(b')=b\), and \(\alpha_2\)
makes it unique.  If \(b'=\operatorname{num}(k)\), then
\(b=S^{\mathcal M}(\operatorname{num}(k))=\operatorname{num}(k+1)\) by Step~0,
a contradiction; so \(b'\) is nonstandard.  By \(\alpha_3\),
\(b'<^{\mathcal M}b\).

Let \(X\) be the set of nonstandard elements, assumed nonempty.  No
\(x\in X\) has the property that no \(y\in X\) satisfies
\(y<^{\mathcal M}x\), since \(x'\in X\) and \(x'<^{\mathcal M}x\).  A least
element \(x\) of \(X\), one with \(x=y\) or \(x<^{\mathcal M}y\) for every
\(y\in X\), would have this property: if \(y<^{\mathcal M}x\) with \(y\in X\),
then \(x=y\) gives \(x<^{\mathcal M}x\) directly and \(x<^{\mathcal M}y\)
gives it by \(\alpha_5\), against \(\alpha_4\).  So \(X\) has no
\(<^{\mathcal M}\)-least element and \(<^{\mathcal M}\) is not a
well-ordering.  Finally, for nonstandard \(b\) put \(b_0=b\) and let
\(b_{k+1}\) be the unique predecessor of \(b_k\).  This is a recursion on
\(\N\) with no choices, since each predecessor is unique and, by what was
just shown, again nonstandard; and \(b_{k+1}<^{\mathcal M}b_k\).

\emph{Clause 5.}  Suppose \(\varphi(x,\bar y)\) and
\(\bar e=(e_1,\ldots,e_l)\) define \(\operatorname{num}[\N]\).  Consider the
induction instance
\[
 \iota_\varphi=\forall\bar y\,\Bigl(\bigl(\varphi(0,\bar y)\land
 \forall x\,(\varphi(x,\bar y)\to\varphi(Sx,\bar y))\bigr)\to
 \forall x\,\varphi(x,\bar y)\Bigr),
\]
where \(\varphi(0,\bar y)\) and \(\varphi(Sx,\bar y)\) are the substitution
instances \(\varphi[0/x]\) and \(\varphi[Sx/x]\).  It is true in
\(\mathcal N_{\mathrm{ar}}\): for parameters \(\bar n'\in\N^l\), the set
\(A=\{x\in\N:\mathcal N_{\mathrm{ar}}\entails\varphi(x,\bar n')\}\) contains
\(0\) when the first conjunct holds and is closed under successor when the
second does (\cref{thm:fol-substitution}), so \(A=\N\) by induction on \(\N\)
in the ambient theory.  Hence \(\mathcal M\entails\iota_\varphi\).  With
\(\bar y\mapsto\bar e\): \(0^{\mathcal M}=\operatorname{num}(0)\) lies in the
defined set, so \(\mathcal M\entails\varphi(0,\bar e)\) by
\cref{thm:fol-substitution}; and if \(\mathcal M\entails\varphi(a,\bar e)\),
then \(a=\operatorname{num}(n)\) for some \(n\), so
\(S^{\mathcal M}(a)=\operatorname{num}(n+1)\) is in the defined set and
\(\mathcal M\entails\varphi(Sx,\bar e)\) at \(x\mapsto a\).  Therefore
\(\mathcal M\entails\forall x\,\varphi(x,\bar e)\), so every element is
standard, contrary to hypothesis.
\end{proof}
\end{annexitem}

\begin{annexitem}{annex:p02-w02-p02-033-direct}{thm:p02-w02-p02-033-nonstandard-direct}{A nonstandard model of true arithmetic from one new constant}
\begin{proof}
\emph{Clause 1.}  Let \(\Delta\subseteq T_1\) be finite, \(m=m_\Delta\), and
let \(\mathcal E\) be the expansion of \(\mathcal N_{\mathrm{ar}}\) with
\(c^{\mathcal E}=m\).  A sentence of
\(\Delta\cap\operatorname{Th}(\mathcal N_{\mathrm{ar}})\) is an
\(\Lang_A\)-sentence true in the \(\Lang_A\)-reduct
\(\mathcal N_{\mathrm{ar}}\) of \(\mathcal E\), hence true in \(\mathcal E\)
by \cref{lem:reduct-invariance}.  A sentence \(\neg\,c=\bar n\) in
\(\Delta\) has \(n\leq m-1<m\); the value of \(\bar n\) in \(\mathcal E\) is
\(n\) (\cref{lem:reduct-invariance} and Step~0 of
\cref{annex:p02-w02-p02-033-numeral}), and the value of \(c\) is \(m\ne n\).
So \(\mathcal E\entails\Delta\).

\emph{Clause 2.}  By clause~1 and \cref{thm:fol-compactness}, clause~1,
\(T_1\) has a model \(\mathcal K\).  The signature \(\Lang_A\cup\{c\}\) is
finite, so \cref{thm:downward-ls} with \(A=\varnothing\) gives
\(\mathcal K_0\preccurlyeq\mathcal K\) with
\(\abs{\dom\mathcal K_0}\leq\aleph_0\).  Elementarity of the inclusion at
the empty tuple gives \(\mathcal K_0\entails T_1\).

\emph{Clause 3.}  By \cref{lem:reduct-invariance},
\(\mathcal M\entails\operatorname{Th}(\mathcal N_{\mathrm{ar}})\), and every
\(\Lang_A\)-term has the same value in \(\mathcal K\) and \(\mathcal M\).
Since \(\mathcal K\entails\neg\,c=\bar n\),
\(c^{\mathcal K}\ne\bar n^{\mathcal M}=\operatorname{num}_{\mathcal M}(n)\)
for every \(n\); so \(c^{\mathcal K}\) is nonstandard.
\Cref{thm:p02-w02-p02-033-numeral-map}, clause~2, gives
\(\mathcal M\entails\bar n<c^{\mathcal K}\); clause~1 makes
\(\operatorname{num}_{\mathcal M}\) elementary, with \(c^{\mathcal K}\)
outside its range.  For an \(\Lang_A\)-sentence \(\sigma\), one of
\(\sigma,\neg\sigma\) lies in \(\operatorname{Th}(\mathcal N_{\mathrm{ar}})\),
so \(\mathcal M\entails\sigma\) iff
\(\mathcal N_{\mathrm{ar}}\entails\sigma\).  Clause~3 of
\cref{thm:p02-w02-p02-033-numeral-map} gives
\(\mathcal M\not\cong\mathcal N_{\mathrm{ar}}\).

\emph{Clause 4.}  Take \(\mathcal K\) with countable domain from clause~2
and its reduct \(\mathcal M\).  By clause~3,
\(\operatorname{num}_{\mathcal M}:\mathcal N_{\mathrm{ar}}\to\mathcal M\) is
an elementary embedding with proper range.
\Cref{thm:p02-w02-p02-033-copy-over-image} gives \(\mathcal M'\) with
\(\mathcal N_{\mathrm{ar}}\preccurlyeq\mathcal M'\),
\(\abs{\dom\mathcal M'}=\abs{\dom\mathcal M}\leq\aleph_0\), and
\(\N\ne\dom\mathcal M'\) because the range of
\(\operatorname{num}_{\mathcal M}\) is proper.
\end{proof}
\end{annexitem}

\begin{annexitem}{annex:p02-w02-p02-033-divisible}{thm:p02-w02-p02-033-divisible}{Divisibility by every prime}
\begin{proof}
\emph{A prime above each number.}  Let \(n\geq1\) and \(u=n!+1\geq2\).  The
least divisor \(q\) of \(u\) with \(q>1\) exists, since \(u\) divides itself,
and is prime, since a divisor of \(q\) strictly between \(1\) and \(q\) would
be a smaller divisor of \(u\) greater than \(1\).  If \(q\leq n\), then
\(q\) divides \(n!\) and \(u\), hence divides \(u-n!=1\), which is
impossible.  So \(q>n\).  Then no \(z\in\N\) has \(q\times z=n\): \(z=0\)
gives \(0\ne n\), and \(z\geq1\) gives \(q\times z\geq q>n\).  Hence
\(\neg\,\bar q\mid\bar n\in\operatorname{Th}(\mathcal N_{\mathrm{ar}})\),
using Steps~0 and~1 of \cref{annex:p02-w02-p02-033-numeral}.

\emph{Clause 1.}  Let \(\Delta\subseteq T_2\) be finite, \(m=m_\Delta\), and
\(\mathcal E\) the expansion of \(\mathcal N_{\mathrm{ar}}\) with
\(c^{\mathcal E}=m\).  Sentences of
\(\Delta\cap\operatorname{Th}(\mathcal N_{\mathrm{ar}})\) hold in
\(\mathcal E\) by \cref{lem:reduct-invariance}.  As a product of primes, or
\(1\), \(m\ne0\), so \(\mathcal E\entails\neg\,c=\bar 0\).  If
\(\bar p\mid c\in\Delta\), then \(p\) is one of the factors of \(m\), and
\(m=p\times z_0\) where \(z_0\) is the product of the remaining factors
(\(1\) if none); the assignment \(z\mapsto z_0\) satisfies
\(\bar p\times z=c\) in \(\mathcal E\), since \(\bar p\) has value \(p\).  So
\(\mathcal E\entails\Delta\).

\emph{Clause 2.}  By clause~1 and \cref{thm:fol-compactness}, clause~1,
\(T_2\) has a model \(\mathcal K\); let \(\mathcal M\) be its reduct and
\(b=c^{\mathcal K}\).  By \cref{lem:reduct-invariance},
\(\mathcal M\entails\operatorname{Th}(\mathcal N_{\mathrm{ar}})\) and
\(\Lang_A\)-terms have the same values in \(\mathcal K\) and \(\mathcal M\).
From \(\mathcal K\entails\neg\,c=\bar 0\) we get
\(b\ne\bar 0^{\mathcal M}\).  From \(\mathcal K\entails\bar p\mid c\) there is
\(e\) with \(\bar p^{\mathcal M}\times^{\mathcal M}e=b\), so
\(\mathcal M\entails\bar p\mid b\).  For \(n\geq1\), the prime \(q>n\) above
gives \(\mathcal M\entails\neg\,\bar q\mid\bar n\).  If \(b\) were
\(\operatorname{num}_{\mathcal M}(n)\), then \(n=0\) is excluded by
\(b\ne\bar 0^{\mathcal M}\), and for \(n\geq1\) the formula
\(\bar q\mid x\) would hold at \(b=\bar n^{\mathcal M}\), so
\(\mathcal M\entails\bar q\mid\bar n\) by \cref{thm:fol-substitution},
contradicting \(\mathcal M\entails\neg\,\bar q\mid\bar n\).  So \(b\) is
nonstandard.
\end{proof}
\end{annexitem}

\begin{annexitem}{annex:p02-w02-p02-033-infinitesimal}{thm:p02-w02-p02-033-infinitesimal}{An elementary extension of the real field with infinitesimals}
\begin{proof}
Write \(\mathcal R\) for \(\mathcal R_{\mathrm{of}}\) and
\(\mathcal R_{\mathbb R}\) for its parameter expansion.  Every
\(\Lang_{\mathrm{of}}\)-sentence true in \(\mathcal R\) is true in
\(\mathcal R_{\mathbb R}\) by \cref{lem:reduct-invariance}, so it belongs to
\(\operatorname{ElDiag}(\mathcal R)\).

\emph{Clause 1.}  Let \(\Delta\subseteq T_3\) be finite, \(r=r_\Delta\), and
\(\mathcal E\) the expansion of \(\mathcal R_{\mathbb R}\) with
\(c^{\mathcal E}=r\).  Sentences of
\(\Delta\cap\operatorname{ElDiag}(\mathcal R)\) are true in
\(\mathcal R_{\mathbb R}\), the
\(\Lang_{\mathrm{of}}(\mathbb R)\)-reduct of \(\mathcal E\), hence in
\(\mathcal E\) by \cref{lem:reduct-invariance}.  If
\(c_a<c\in\Delta\), then \(a\leq\max\{a':c_{a'}<c\in\Delta\}<r\), so
\(\mathcal E\entails c_a<c\).

\emph{Clause 2.}  Compactness (\cref{thm:fol-compactness}, clause~1) gives
\(\mathcal K\entails T_3\).  Let \(\mathcal K_0\) be the
\(\Lang_{\mathrm{of}}(\mathbb R)\)-reduct of \(\mathcal K\); it satisfies
\(\operatorname{ElDiag}(\mathcal R)\) by \cref{lem:reduct-invariance}, and
\(\mathcal M\) is its \(\Lang_{\mathrm{of}}\)-reduct.  Let
\(\varphi(x_1,\ldots,x_k)\) be an \(\Lang_{\mathrm{of}}\)-formula and
\(a_1,\ldots,a_k\in\mathbb R\).  By \cref{thm:fol-substitution}, as in
Step~1 of \cref{annex:p02-w02-p02-033-numeral},
\(\mathcal R\entails\varphi(a_1,\ldots,a_k)\) iff
\(\varphi(c_{a_1},\ldots,c_{a_k})\in\operatorname{ElDiag}(\mathcal R)\), and
otherwise its negation belongs to \(\operatorname{ElDiag}(\mathcal R)\).
Since \(\mathcal K_0\) satisfies the one that belongs, and by
\cref{thm:fol-substitution} and \cref{lem:reduct-invariance}
\(\mathcal K_0\entails\varphi(c_{a_1},\ldots,c_{a_k})\) iff
\(\mathcal M\entails\varphi(j(a_1),\ldots,j(a_k))\), the map \(j\) is
elementary.  (This is the second test of \cref{thm:diagram-tests}.)  For
\(a\in\mathbb R\), \(\mathcal K\entails c_a<c\) gives
\(\mathcal M\entails j(a)<c^{\mathcal K}\).

By induction on \(n\), \(\underline n^{\mathcal R}=n\).  By
\cref{thm:p02-w02-p02-033-elementary-maps-embed}, \(j\) is an embedding,
hence a homomorphism, so \cref{lem:term-preservation} gives
\(\underline n^{\mathcal M}=j(n)\).  The sentences
\(\beta_1=\forall x\,\neg(x<x)\) and
\(\beta_2=\forall x\,\forall y\,\forall w\,(x<y\land y<w\to x<w)\) are true in
\(\mathcal R\), hence in \(\mathcal M\)
(\cref{thm:p02-w02-p02-033-elementary-maps-embed}).  If
\(\mathcal M\entails c^{\mathcal K}<\underline n\), then with
\(\mathcal M\entails j(n)<c^{\mathcal K}\) and \(\beta_2\) we would get
\(\mathcal M\entails c^{\mathcal K}<c^{\mathcal K}\), against \(\beta_1\).
So \(x=c^{\mathcal K}\) has no \(n\) with
\(\mathcal M\entails x<\underline n\), and \(\mathcal M\) is not
Archimedean.

\emph{Clause 3.}  The \(\Lang_{\mathrm{of}}\)-sentences
\[
 \gamma_1=\forall x\,\bigl(0<x\to\exists y\,(x\cdot y=1)\bigr),\qquad
 \gamma_2=\forall x\,\forall y\,\bigl((0<x\land x\cdot y=1)\to0<y\bigr)
\]
are true in \(\mathcal R\): a positive real has a multiplicative inverse,
and the inverse of a positive real is positive.  For each real \(a>0\), the
\(\Lang_{\mathrm{of}}(\mathbb R)\)-sentence
\[
 \delta_a=\forall x\,\forall y\,\bigl((c_{1/a}<x\land x\cdot y=1)\to
 y<c_a\bigr)
\]
is true in \(\mathcal R_{\mathbb R}\): if \(x>1/a\) and \(xy=1\), then
\(x>0\), \(1<ax\), and \(y=1/x<a\).  So
\(\gamma_1,\gamma_2,\delta_a\in\operatorname{ElDiag}(\mathcal R)\subseteq T_3\),
and \(\mathcal K\) satisfies them.  The sentence \(c_0=0\) is in
\(\operatorname{ElDiag}(\mathcal R)\) and \(c_0<c\) is in \(T_3\), so
\(\mathcal K\entails0<c\).  By \(\gamma_1\) there is \(d\) with
\(c^{\mathcal K}\cdot^{\mathcal K}d=1^{\mathcal K}\), and by \(\gamma_2\),
\(\mathcal K\entails0<d\) at \(d\).  For real \(a>0\),
\(\mathcal K\entails c_{1/a}<c\), so \(\delta_a\) with
\(x\mapsto c^{\mathcal K}\) and \(y\mapsto d\) gives
\(d<^{\mathcal K}c_a^{\mathcal K}\).  Passing to the reduct,
\(\mathcal M\entails0<d\) and \(\mathcal M\entails d<j(a)\).  Suppose
\(d=j(b)\).  As \(j\) preserves constants, \(0^{\mathcal M}=j(0)\), so
\(\mathcal M\entails j(0)<j(b)\) and elementarity gives \(0<b\).  Then
\(a=b\) yields \(\mathcal M\entails d<d\), against \(\beta_1\).  So
\(d\notin j[\mathbb R]\).

\emph{Clause 4.}  \Cref{thm:p02-w02-p02-033-copy-over-image}, applied to
\(j\), gives \(\mathcal M'\) with \(\mathcal R\preccurlyeq\mathcal M'\) and an
isomorphism \(\pi:\mathcal M'\to\mathcal M\) with \(\pi(a)=j(a)\) for
\(a\in\mathbb R\).  Put \(\varepsilon=\pi^{-1}(d)\).  If
\(\varepsilon=a\in\mathbb R\), then \(d=\pi(a)=j(a)\), which clause~3
excludes.  Since \(\mathcal R\subseteq\mathcal M'\), the inclusion preserves
constants and \(0^{\mathcal M'}\) is the real number \(0\); as \(\pi\)
preserves constants and preserves and reflects \(<\),
\(\mathcal M'\entails0<\varepsilon\) iff \(\mathcal M\entails0<d\), and
\(\mathcal M'\entails\varepsilon<a\) iff \(\mathcal M\entails d<j(a)\), for
each real \(a>0\).  Both right-hand sides hold by clause~3.
\end{proof}
\end{annexitem}

\begin{annexitem}{annex:p02-w02-p02-033-fails}{thm:p02-w02-p02-033-fs-fails}{Where finite satisfiability fails}
\begin{proof}
\emph{Clause 1.}  Let
\(\sigma_{\mathcal F}=\forall x\,(x=c_{a_1}\lor\cdots\lor x=c_{a_k})\).  Every
element of \(\dom\mathcal F\) is some \(a_i\), the value of \(c_{a_i}\) in
the parameter expansion, so \(\sigma_{\mathcal F}\in
\operatorname{ElDiag}(\mathcal F)\).  Let
\(\Delta_0=\{\sigma_{\mathcal F}\}\cup\{\neg\,c=c_{a_i}:1\leq i\leq k\}\).  If
\(\mathcal K\entails\Delta_0\), then \(\sigma_{\mathcal F}\) at
\(x\mapsto c^{\mathcal K}\) gives \(c^{\mathcal K}=c_{a_i}^{\mathcal K}\) for
some \(i\), against \(\neg\,c=c_{a_i}\).  So \(\Delta_0\) has no model, and
neither has any set containing it.

\emph{Clause 2.}  The finite-satisfiability claims are clause~1 of
\cref{thm:p02-w02-p02-033-nonstandard-direct},
\cref{thm:p02-w02-p02-033-divisible} and
\cref{thm:p02-w02-p02-033-infinitesimal}.  An expansion of
\(\mathcal N_{\mathrm{ar}}\) interprets \(c\) as some \(m\in\N\).  It
falsifies \(\neg\,c=\bar m\in T_1\), since \(\bar m\) has value \(m\).  For
\(T_2\): if \(m=0\) it falsifies \(\neg\,c=\bar 0\); if \(m\geq1\), the
prime \(q>m\) of \cref{annex:p02-w02-p02-033-divisible} satisfies
\(q\times z\ne m\) for every \(z\in\N\), so the expansion falsifies
\(\bar q\mid c\in T_2\).  An expansion of \(\mathcal R_{\mathbb R}\)
interprets \(c\) as some \(r\in\mathbb R\) and falsifies \(c_r<c\in T_3\),
since \(r<r\) fails.
\end{proof}
\end{annexitem}

\begin{annexitem}{annex:p02-w02-p02-033-notax}{thm:p02-w02-p02-033-not-axiomatizable}{Three properties with no first-order axiomatization}
\begin{proof}
\emph{Clause 1.}  Suppose \(\Sigma\) had exactly the structures isomorphic
to \(\mathcal N_{\mathrm{ar}}\) as models.  Then
\(\mathcal N_{\mathrm{ar}}\entails\Sigma\), so
\(\Sigma\subseteq\operatorname{Th}(\mathcal N_{\mathrm{ar}})\).  By
\cref{thm:p02-w02-p02-033-nonstandard-direct}, clauses~2 and~3, there is an
\(\Lang_A\)-structure \(\mathcal M\) satisfying
\(\operatorname{Th}(\mathcal N_{\mathrm{ar}})\), hence \(\Sigma\), with
\(\mathcal M\not\cong\mathcal N_{\mathrm{ar}}\), a contradiction.

\emph{Clause 2.}  With \(\Sigma\subseteq\operatorname{Th}(\mathcal N_{\mathrm{ar}})\),
the same \(\mathcal M\) satisfies \(\Sigma\) and has the nonstandard element
\(c^{\mathcal K}\); by clause~4 of \cref{thm:p02-w02-p02-033-numeral-map},
\(<^{\mathcal M}\) is not a well-ordering.

\emph{Clause 3.}  Let \(\Sigma\) be true in \(\mathcal R_{\mathrm{of}}\).
Take \(\mathcal M\) and \(j\) from clause~2 of
\cref{thm:p02-w02-p02-033-infinitesimal}.  By
\cref{thm:p02-w02-p02-033-elementary-maps-embed}, \(\mathcal M\) satisfies
every sentence true in \(\mathcal R_{\mathrm{of}}\), so
\(\mathcal M\entails\Sigma\); and \(\mathcal M\) is not Archimedean.
\end{proof}
\end{annexitem}

\begin{annexitem}{annex:p02-w02-p02-033-cantor}{thm:p02-w02-p02-033-cantor-internal}{The theory proves that an uncountable set exists}
\begin{proof}
Let \(\mathcal M\models T\).  Write \(\operatorname{ext}\) for
\(\operatorname{ext}_{\mathcal M}\) and \(b\in^{\mathcal M}a\) for
\(\mathcal M\entails b\in a\).  By the satisfaction clauses, a bounded
quantifier \(\forall x\in a\) (respectively \(\exists x\in a\)) holds in
\(\mathcal M\) iff the matrix holds for every (respectively some)
\(x\in\operatorname{ext}(a)\).  Extensionality in \(\mathcal M\) says:
if \(\operatorname{ext}(a)=\operatorname{ext}(a')\), then \(a=a'\).

\emph{Clause 1, existence of \(\omega^{\mathcal M}\).}  By the axiom of
infinity there is \(I\in\dom\mathcal M\) with
\(\mathcal M\entails\mathrm{Ind}(I)\).  The separation instance for the
formula \(\forall v\,(\mathrm{Ind}(v)\to x\in v)\), applied to \(I\), gives
\(w\) with
\[
 \operatorname{ext}(w)=\{x\in\operatorname{ext}(I):
 x\in^{\mathcal M}v\text{ for every }v\text{ with }
 \mathcal M\entails\mathrm{Ind}(v)\}.
\]
Leastness: every \(x\in\operatorname{ext}(w)\) lies in every \(v\) with
\(\mathcal M\entails\mathrm{Ind}(v)\), which is the second conjunct of
\(\mathrm{Om}(w)\).  Inductivity: \(I\) contains some \(e\) with
\(\operatorname{ext}(e)=\varnothing\).  Every \(v\) with
\(\mathcal M\entails\mathrm{Ind}(v)\) contains some \(e'\) with
\(\operatorname{ext}(e')=\varnothing\), and extensionality gives \(e'=e\);
so \(e\in^{\mathcal M}v\) for all such \(v\), and \(e\in\operatorname{ext}(w)\).
Next let \(x\in\operatorname{ext}(w)\).  As \(x\in^{\mathcal M}I\), there is
\(y\in^{\mathcal M}I\) with
\(\operatorname{ext}(y)=\operatorname{ext}(x)\cup\{x\}\).  For every \(v\)
with \(\mathcal M\entails\mathrm{Ind}(v)\) we have \(x\in^{\mathcal M}v\), so
some \(y'\in^{\mathcal M}v\) has
\(\operatorname{ext}(y')=\operatorname{ext}(x)\cup\{x\}\); extensionality
gives \(y'=y\).  Hence \(y\in\operatorname{ext}(w)\), and
\(\mathcal M\entails\mathrm{Ind}(w)\).  So
\(\mathcal M\entails\mathrm{Om}(w)\).

\emph{Clause 1, uniqueness and \(\mathcal P(\omega)^{\mathcal M}\).}  If
\(\mathcal M\entails\mathrm{Om}(w)\land\mathrm{Om}(w')\), leastness of each
applied to the inductivity of the other gives
\(\operatorname{ext}(w)\subseteq\operatorname{ext}(w')\) and conversely, so
\(w=w'\).  The power-set axiom of \cref{sec:zf-zfc-01}, with \(x\subseteq a\)
unabbreviated, is \(\forall a\,\exists p\,\mathrm{Pw}(p,a)\); it gives
\(p\) with \(\mathcal M\entails\mathrm{Pw}(p,w)\), and two such \(p\) have
the same external extension, hence are equal.

\emph{Clause 2.}  Write \(w=\omega^{\mathcal M}\) and
\(p=\mathcal P(\omega)^{\mathcal M}\), and suppose, for a contradiction,
that \(g\in\dom\mathcal M\) satisfies
\(\mathcal M\entails\mathrm{Inj}(g,p,w)\).  The separation instance for the
formula
\[
 \exists x\,\bigl(x\in q\land\mathrm{In}(x,n,g)\land\neg(n\in x)\bigr),
\]
with free variable \(n\) and parameters \(q,g\), applied to \(w\) with
\(q:=p\), gives \(d\in\dom\mathcal M\) with
\[
\begin{aligned}
 \operatorname{ext}(d)=\{n\in\operatorname{ext}(w):{}&\text{some }
 x\in\operatorname{ext}(p)\text{ has }\mathcal M\entails\mathrm{In}(x,n,g)\\
 &\text{and }n\notin\operatorname{ext}(x)\}.
\end{aligned}
\]
Since \(\operatorname{ext}(d)\subseteq\operatorname{ext}(w)\), the formula
\(\forall t\,(t\in d\to t\in w)\) holds, and \(\mathrm{Pw}(p,w)\) gives
\(d\in\operatorname{ext}(p)\).  The first conjunct of
\(\mathrm{Inj}(g,p,w)\) gives \(k\in\operatorname{ext}(w)\) with
\(\mathcal M\entails\mathrm{In}(d,k,g)\).
If \(k\in\operatorname{ext}(d)\), some \(x\in\operatorname{ext}(p)\) has
\(\mathcal M\entails\mathrm{In}(x,k,g)\) and
\(k\notin\operatorname{ext}(x)\); the third conjunct of
\(\mathrm{Inj}(g,p,w)\), applied to \(x,d,k\), gives \(x=d\), so
\(k\notin\operatorname{ext}(d)\), a contradiction.
If \(k\notin\operatorname{ext}(d)\), then \(x:=d\) witnesses the defining
condition of \(d\) for \(n:=k\): \(d\in\operatorname{ext}(p)\),
\(\mathcal M\entails\mathrm{In}(d,k,g)\) and \(k\notin\operatorname{ext}(d)\);
so \(k\in\operatorname{ext}(d)\), a contradiction.  Hence no such \(g\)
exists, and \(w,p\) witness \(\sigma_{\mathrm{unc}}\) in \(\mathcal M\).
As \(\mathcal M\) was an arbitrary model of \(T\),
\(T\entails\sigma_{\mathrm{unc}}\).

\emph{Clause 3.}  Assuming the Axiom of Choice,
\cref{thm:fol-completeness} turns \(T\entails\sigma_{\mathrm{unc}}\) into
\(T\derives_{\mathsf{ND}_{\mathsf{FOL}}}\sigma_{\mathrm{unc}}\), and
\cref{thm:fol-soundness} returns \(T\entails\sigma_{\mathrm{unc}}\) from that
derivation.  Finally, in any model \(\mathcal M\), the witnesses \(w,p\) of
\(\sigma_{\mathrm{unc}}\) are \(\omega^{\mathcal M}\) and
\(\mathcal P(\omega)^{\mathcal M}\) by clause~1.  If
\(\mathcal M\entails\mathrm{Ctbl}(p)\), its witness \(w'\) satisfies
\(\mathrm{Om}\), so \(w'=w\) by clause~1, and some \(g\) satisfies
\(\mathrm{Inj}(g,p,w)\), contrary to clause~2.  Hence
\(\mathcal M\entails\neg\mathrm{Ctbl}(p)\) and
\(\mathcal M\entails\exists a\,\neg\mathrm{Ctbl}(a)\).
\end{proof}
\end{annexitem}

\begin{annexitem}{annex:p02-w02-p02-033-skolem}{thm:p02-w02-p02-033-skolem-resolution}{Skolem's paradox: internally uncountable, externally countable}
\begin{proof}
\emph{Clause 1.}  Let \(\mathcal M_0\) be a model of \(T\).
\Cref{thm:downward-ls} with \(A=\varnothing\) gives
\(\mathcal M\preccurlyeq\mathcal M_0\) with
\(\abs{\dom\mathcal M}\leq\max(0,\abs{\Lang_{\in}},\aleph_0)=\aleph_0\), so
\(\dom\mathcal M\) is countable.  Elementarity at the empty tuple makes
\(\mathcal M\) satisfy every sentence of \(T\).

\emph{Clause 2.}  \(\mathcal M\entails\neg\mathrm{Ctbl}(p)\) is the last
part of the proof of clause~3 of
\cref{thm:p02-w02-p02-033-cantor-internal}, which used only clauses 1 and~2
of that theorem.  As \(\dom\mathcal M\) is countable there is an injection
\(\iota:\dom\mathcal M\to\N\), and its restriction to the subset
\(\operatorname{ext}_{\mathcal M}(p)\) is an injection into \(\N\).

\emph{Clause 3.}  Write \(P=\operatorname{ext}_{\mathcal M}(p)\) and
\(W=\operatorname{ext}_{\mathcal M}(w)\).  By the satisfaction clauses for
bounded quantifiers, the three conjuncts of \(\mathrm{Inj}(g,p,w)\) hold in
\(\mathcal M\) iff, respectively, every \(x\in P\) has some \(n\in W\) with
\((x,n)\in R_g\); every \(x\in P\) has at most one such \(n\in W\); and no
two distinct \(x,x'\in P\) share an \(n\in W\) with \((x,n),(x',n)\in R_g\).
These three conditions together say exactly that
\(R_g\cap(P\times W)\) is an injective function from \(P\) into \(W\).
Clause~2 of \cref{thm:p02-w02-p02-033-cantor-internal} shows that no
\(g\in\dom\mathcal M\) satisfies \(\mathrm{Inj}(g,p,w)\) in \(\mathcal M\),
so no \(g\in\dom\mathcal M\) has this property.

\emph{Clause 4.}  Define \(e_0\) as the unique element with
\(\operatorname{ext}_{\mathcal M}(e_0)=\varnothing\), which lies in \(W\)
because \(\mathcal M\entails\mathrm{Ind}(w)\) and extensionality makes the
empty element unique; and, given \(e_k\in W\), define \(e_{k+1}\) as the
unique element of \(W\) with
\(\operatorname{ext}_{\mathcal M}(e_{k+1})=
\operatorname{ext}_{\mathcal M}(e_k)\cup\{e_k\}\), which exists by
\(\mathrm{Ind}(w)\) and is unique by extensionality.  This is a recursion on
\(\N\) in the ambient theory with no choices.  By induction on \(k\),
\(\operatorname{ext}_{\mathcal M}(e_k)=\{e_0,\ldots,e_{k-1}\}\) and
\(e_0,\ldots,e_k\) are pairwise distinct.  For \(k=0\) both are immediate.
Assume both for \(k\).  Then
\(\operatorname{ext}_{\mathcal M}(e_{k+1})=\{e_0,\ldots,e_k\}\) has
\(k+1\) elements, while for \(i\leq k\) the set
\(\operatorname{ext}_{\mathcal M}(e_i)=\{e_0,\ldots,e_{i-1}\}\) has
\(i\leq k\) elements, so \(e_{k+1}\ne e_i\).  Hence \(k\mapsto e_k\) is an
injection \(\N\to W\), and \(b\mapsto e_{\iota(b)}\) is an injective
function from \(P\) into \(W\) in the ambient set theory.  By clause~3 it is
\(R_g\cap(P\times W)\) for no \(g\in\dom\mathcal M\).
\end{proof}
\end{annexitem}
